\documentclass[reqno]{amsart}
\newcommand{\problemnumber}{12480}

\usepackage{enumerate}
\usepackage{color}
% \usepackage[top=3cm, bottom=3cm, left=3cm, right=3cm]{geometry}

\theoremstyle{definition}
\newtheorem*{solution}{Solution}
\newtheorem*{remark}{Remark}

\newcommand*\diff{\mathop{}\!\mathrm{d}}
\newcommand{\abs}[1]{\left\lvert #1 \right\rvert}
\DeclareMathOperator{\esp}{e}
\usepackage{upgreek}

\newtheorem{claim}{Claim}

\newcommand{\fulges}[1]{{\color{cyan}#1}}


%% shortcuts
\newcommand{\bbZ}{\mathbb{Z}}
\newcommand{\bbN}{\mathbb{N}}
\newcommand{\bbR}{\mathbb{R}}
\newcommand{\vvirg}{, \ldots ,}

\title{Solution to AMM Problem \problemnumber}

\author[I.Bomze]{Immanuel Bomze}
\email{}
\author[T. Mannelli Mazzoli]{Tommaso Mannelli Mazzoli}
\email{tommaso.mazzoli@tuwien.ac.at}

\setlength{\parskip}{.3cm}
\begin{document}
\maketitle

\subsection*{Problem \problemnumber}\textit{Proposed by F. Stanescu (Romania)} \\

Let $f$ and $g$ be two nonnegative functions on $[0,1]$ with $f(0)=g(0)=1$,
and let $h\colon [0,1] \to \mathbb{R}$ be nondecreasing. Suppose that $f$ is convex and $g$ is concave. Prove
\[
\int_0^1 g(x) \diff x \int_0^1 f(x)h(x) \diff x \ge \int_0^1 f(x) \diff x 
\int_0^1 g(x)h(x) \diff x 
\]

\begin{solution}
\textit{\noindent
We show a slightly more general statement: 
let $f$ and $g$ be two nonnegative functions on $[0,1]$ with $f(0)=g(0)=\alpha\ge0
$,
and let $h\colon [0,1] \to \mathbb{R}$ be nondecreasing. Suppose that $f$ is convex and $g$ is concave. Then
\[
\int_0^1 g(x) \diff x \int_0^1 f(x)h(x) \diff x \ge \int_0^1 f(x) \diff x 
\int_0^1 g(x)h(x) \diff x 
\]
}


First, let us consider the case in which $\int_0^1 f(x) \diff x = 0$.
Since $f$ is convex in a compact, it is continuous. Therefore, since $f$ is nonnegative, the condition $\int_0^1 f(x) \diff x = 0$ implies $f \equiv 0$ in $[0,1]$. The desired relation follows with equality.
The same argument applies to $g$.

So let us assume that $\int_0^1 f(x) \diff x > 0 < \int_0^1 g(x) \diff x$. Then, we can divide both sides by $
\int_0^1 f(x) \diff x \cdot \int_0^1 g(x) \diff x $. We can assume then, without loss of generality, that
\[
    \int_0^1 f(x) \diff x = \int_0^1 g(x) \diff x = 1\, .
\]
The problem can be then rewritten as:
\[
    \int_0^1 h(x) f(x) \diff x \ge \int_0^1 h(x) g(x) \diff x\]
    or equivalently 
    \[\int_0^1 \Bigl(f(x)-g(x)\Bigr)h(x) \diff x \ge 0 \, .
\]
Note that
\begin{itemize}
    \item $f-g$ is a convex function, since $-g$ is convex;
    \item $f(0)-g(0)=\alpha-\alpha=0$, by hypothesis;
    \item  $\int_0^1 (f(x)-g(x)) \diff x = 0$, because we assumed $\int_0^1 f(x) \diff x = 1 = \int_0^1 g(x) \diff x $.
\end{itemize}
Hence, there must be a point $t_0\in [0,1)$ such that $f(x)-g(x) \le 0$ for $x\in (0,t_0)$ and $f(x)-g(x) \ge 0$ for $x\in (t_0,1)$.

Therefore
\[
\begin{split}
    \int_0^1 \Bigl(f(x)-g(x)\Bigr)h(x) \diff x 
    &= \int_0^{t_0} \Bigl(f(x)-g(x)\Bigr)h(x) \diff x  + \int_{t_0}^1 \Bigl(f(x)-g(x)\Bigr)h(x) \diff x  \\
    &\ge h(t_0)\int_0^{t_0} (f(x)-g(x))\diff x +h(t_0)\int_{t_0}^1 (f(x)-g(x))\diff x \\
    &= h(t_0)\int_0^{1} (f(x)-g(x))\diff x = 0       
\end{split},
\]
where we used the fact that $f(x)-g(x) \le 0$ in $[0,t_0]$.

\end{solution}


\end{document}